% GENERATED FILE. Edit canonical lessons, then run npm run generate:books.
% canonical-source-hash: fnv1a64:cc09edd5508f3565
% canonical-lessons: ZH-C174-gua, ZH-C174-nong, ZH-C174-tai, ZH-C174-tizi, ZH-C174-dengpao

\chapter{Jobs Around the House}
\label{ch:zh-jobs-around-the-house}
\hlchaptermodality{174}

\begin{chapteropening}
\textbf{By the end of this chapter:} \emph{I can hang things up, sort things out, lift them together, and talk about a ladder and a light bulb.}

Complete the last lesson of chapter 174: I can hang things up, sort things out, lift them together, and talk about a ladder and a light bulb.
\end{chapteropening}

\section[guà]{guà — to hang; to hang up}

\label{lesson:ZH-C174-gua}



\begin{quote}
Before the new one: say the Mandarin for a get-together, then the Mandarin for to pay for everyone.
\end{quote}

\subsection*{You'll want to know: guà}
\emph{guà} — \textquotedblleft{}to hang; to hang up\textquotedblright{}.

\begin{quote}
\emph{mén hòubian kěyǐ guà dàyī} — You can hang your coat behind the door.
\end{quote}

Its characters come later; for now it is a word for the ear.

\begin{cousinweb}
\emph{guà} is inherited from Middle Chinese \emph{kweaH} (Guangyun, in Baxter's transcription), glossed 'to hang', under a popular way of writing it.
\end{cousinweb}

\begin{grammarlens}[title={What you've built}]
One more word for saying what I do.
\end{grammarlens}

\subsection*{Your turn}
\begin{itemize}
\raggedright
  \item \emph{Say it:} \emph{guà}
  \item \emph{Say it:} \emph{guà}, once more
  \item \emph{Say it:} \emph{mén hòubian kěyǐ guà dàyī} — You can hang your coat behind the door.
  \item \emph{Recall:} say the Mandarin for a get-together, then the Mandarin for to pay for everyone, then say \emph{guà} again
\end{itemize}

\begin{tcolorbox}[breakable,colback=teal!4,colframe=teal!35!black,title={Before you move on}]
What does guà mean? (\textquotedblleft{}To hang; to hang up\textquotedblright{}.) Say it once more.
\end{tcolorbox}
\section[nòng]{nòng — to sort something out; to do; to handle}

\label{lesson:ZH-C174-nong}



\begin{quote}
Before the new one: say the Mandarin for to pay for everyone, then the Mandarin for to hang something up.
\end{quote}

\subsection*{You'll want to know: nòng}
\emph{nòng} — \textquotedblleft{}to sort something out; to do; to handle\textquotedblright{}.

\begin{quote}
\emph{diànnǎo huài le nǐ huì nòng ma} — The computer's broken: do you know how to sort it out?
\end{quote}

Its characters come later; for now it is a word for the ear.

\begin{cousinweb}
\emph{nòng} is inherited from Middle Chinese \emph{luwngH} (Guangyun, in Baxter's transcription), glossed 'to play with'. Its first sound would regularly give lòng; how it came to be said with n is not settled.
\end{cousinweb}

\begin{grammarlens}[title={What you've built}]
One more word for saying what I do.
\end{grammarlens}

\subsection*{Your turn}
\begin{itemize}
\raggedright
  \item \emph{Say it:} \emph{nòng}
  \item \emph{Say it:} \emph{nòng}, once more
  \item \emph{Say it:} \emph{diànnǎo huài le nǐ huì nòng ma} — The computer's broken: do you know how to sort it out?
  \item \emph{Recall:} say the Mandarin for to pay for everyone, then the Mandarin for to hang something up, then say \emph{nòng} again
\end{itemize}

\begin{tcolorbox}[breakable,colback=teal!4,colframe=teal!35!black,title={Before you move on}]
What does nòng mean? (\textquotedblleft{}To sort something out; to do; to handle\textquotedblright{}.) Say it once more.
\end{tcolorbox}
\section[tái]{tái — to lift; to carry (between two people)}

\label{lesson:ZH-C174-tai}



\begin{quote}
Before the new one: say the Mandarin for to hang something up, then the Mandarin for to sort something out.
\end{quote}

\subsection*{You'll want to know: tái}
\emph{tái} — \textquotedblleft{}to lift; to carry (between two people)\textquotedblright{}.

\begin{quote}
\emph{zhège zhuōzi hěn zhòng wǒmen yìqǐ tái} — This table is heavy: let's lift it together.
\end{quote}

Its characters come later; for now it is a word for the ear.

\begin{cousinweb}
\emph{tái} is inherited from Middle Chinese \emph{doj} (Guangyun, in Baxter's transcription), written there another way and glossed 'to lift up'.
\end{cousinweb}

\begin{grammarlens}[title={What you've built}]
One more word for saying what I do.
\end{grammarlens}

\subsection*{Your turn}
\begin{itemize}
\raggedright
  \item \emph{Say it:} \emph{tái}
  \item \emph{Say it:} \emph{tái}, once more
  \item \emph{Say it:} \emph{zhège zhuōzi hěn zhòng wǒmen yìqǐ tái} — This table is heavy: let's lift it together.
  \item \emph{Recall:} say the Mandarin for to hang something up, then the Mandarin for to sort something out, then say \emph{tái} again
\end{itemize}

\begin{tcolorbox}[breakable,colback=teal!4,colframe=teal!35!black,title={Before you move on}]
What does tái mean? (\textquotedblleft{}To lift; to carry (between two people)\textquotedblright{}.) Say it once more.
\end{tcolorbox}
\section[tīzi]{tīzi — ladder}

\label{lesson:ZH-C174-tizi}



\begin{quote}
Before the new one: say the Mandarin for to sort something out, then the Mandarin for to lift together.
\end{quote}

\subsection*{You'll want to know: tīzi}
\emph{tīzi} — \textquotedblleft{}ladder\textquotedblright{}.

\begin{quote}
\emph{wǒmen yìqǐ tái tīzi} — Let's carry the ladder together.
\end{quote}

Its characters come later; for now it is a word for the ear.

\begin{cousinweb}
\emph{tīzi} is built from \emph{tī} and the noun ending \zh{子}, both inherited from Middle Chinese, \emph{thej} and \emph{tsiX} (Guangyun, in Baxter's transcription), which glosses the first 'wooden steps'. In the word the second loses its tone.
\end{cousinweb}

\begin{grammarlens}[title={What you've built}]
One more word for saying what I do.
\end{grammarlens}

\subsection*{Your turn}
\begin{itemize}
\raggedright
  \item \emph{Say it:} \emph{tīzi}
  \item \emph{Say it:} \emph{tīzi}, once more
  \item \emph{Say it:} \emph{wǒmen yìqǐ tái tīzi} — Let's carry the ladder together.
  \item \emph{Recall:} say the Mandarin for to sort something out, then the Mandarin for to lift together, then say \emph{tīzi} again
\end{itemize}

\begin{tcolorbox}[breakable,colback=teal!4,colframe=teal!35!black,title={Before you move on}]
What does tīzi mean? (\textquotedblleft{}Ladder\textquotedblright{}.) Say it once more.
\end{tcolorbox}
\section[dēngpào]{dēngpào — light bulb}

\label{lesson:ZH-C174-dengpao}



\begin{quote}
Before the new one: say the Mandarin for to lift together, then the Mandarin for a ladder.
\end{quote}

\subsection*{You'll want to know: dēngpào}
\emph{dēngpào} — \textquotedblleft{}light bulb\textquotedblright{}.

\begin{quote}
\emph{dēngpào huài le wǒmen yào huàn yí ge} — The bulb's gone, so we need to change it.
\end{quote}

Its characters come later; for now it is a word for the ear.

\begin{cousinweb}
\emph{dēngpào} is built from \emph{dēng} 'lamp' and \emph{pào} 'bubble', both inherited from Middle Chinese, \emph{tong} and \emph{phaew} (Guangyun, in Baxter's transcription), which glosses the second 'bubbles on water': the lamp's glass bubble. \emph{Phaew} would regularly give a first tone, so how \emph{pào} came about is not settled.
\end{cousinweb}

\begin{grammarlens}[title={What you've built}]
One more word for saying what I do.
\end{grammarlens}

\subsection*{Your turn}
\begin{itemize}
\raggedright
  \item \emph{Say it:} \emph{dēngpào}
  \item \emph{Say it:} \emph{dēngpào}, once more
  \item \emph{Say it:} \emph{dēngpào huài le wǒmen yào huàn yí ge} — The bulb's gone, so we need to change it.
  \item \emph{Recall:} say the Mandarin for to lift together, then the Mandarin for a ladder, then say \emph{dēngpào} again
\end{itemize}

\begin{tcolorbox}[breakable,colback=teal!4,colframe=teal!35!black,title={Before you move on}]
What does dēngpào mean? (\textquotedblleft{}Light bulb\textquotedblright{}.) Say it once more.
\end{tcolorbox}
